\subsection{紧凸集的极点}

定义 1. --- 设 A 是仿射空间 E 中的一个凸集。若不存在既含于 A 又包含 x 的开线段，则称点 \(x \in A\) 为 A 的一个极点。

换言之，诸关系式 \(x=\lambda y+(1-\lambda)z\)、\(y\in A\)、\(z\in A\)、\(y\neq z\) 与 \(0\leqslant\lambda\leqslant1\) 蕴含 \(\lambda=0\) 或 \(\lambda=1\)（于是 x=y 或 x=z）。由此推出，x 不可能是 A 中带正质量的 n 个点 \(x_{i}\) 所成之集的重心，除非 x 是诸 \(x_{i}\) 之一；因为 n=2 时这正是定义；对任意的 n，则对 n 作归纳论证：由于 x 是 \(x_{1}\) 与 \(y_{1}\)——后者是诸 \(x_{i}\)（\(2\leqslant i\leqslant n\)）的重心——两者的重心，因此 x 与 \(x_{1}\) 或 \(y_{1}\) 重合，而在第二种情形只需应用归纳假设。

说 \(x\) 是 \(A\) 的极点，也等价于说 \(A - \{x\}\) 是凸的。

例. --- 1) 在空间 \(\mathbf{R}^n\) 中，球面 \(\mathbf{S}_{n-1}\) 的所有点都是闭球 \(\mathbf{B}_n\) 的极点。因为，若 \(\sum_{i} y_i^2 \leqslant 1\)、\(\sum_{i} z_i^2 \leqslant 1\) 且 \(0 < \lambda < 1\)，则关系式

\[
\lambda^ {2} \sum_ {i} y _ {i} ^ {2} + (1 - \lambda) ^ {2} \sum_ {i} z _ {i} ^ {2} + 2 \lambda (1 - \lambda) \sum_ {i} y _ {i} z _ {i} = 1 = (\lambda + (1 - \lambda)) ^ {2}
\]

仅在下述条件成立时才可能：

\[
\sum_ {i} y _ {i} ^ {2} = \sum_ {i} z _ {i} ^ {2} = \sum_ {i} y _ {i} z _ {i} = 1.
\]

但这蕴含 \(\sum_{i}(y_i - z_i)^2 = 0\)，从而 \(y_i = z_i\) 对所有 \(i\) 成立，这就证明了我们的断言。

\begin{enumerate}
\def\labelenumi{\arabic{enumi})}
\setcounter{enumi}{1}
\tightlist
\item
  在实数有界序列组成的赋范空间 \(\mathcal{B}(\mathbf{N})\)（I, p.~4）中，单位球的极点是点 \(x = (\xi_n)\)，它满足 \(|\xi_n| = 1\)（对所有 \(n\)）。因为，假设 \(|\xi_n| \leqslant 1\) 对所有 \(n\) 成立，而 \(|\xi_p| < 1\) 对某个指标 \(p\) 成立。这时我们可以写
\end{enumerate}

\[
x = \frac {1 + \xi_ {p}}{2} y + \frac {1 - \xi_ {p}}{2} z
\]

其中 \(y\)（相应地，\(z\)）是这样的点：它的所有坐标都等于 \(x\) 的同下标坐标，唯独下标 \(p\) 的坐标等于 1（相应地，- 1）。这表明 \(x\) 不是极点，因为我们有 \(\| y \| \leqslant 1\) 与 \(\| z \| \leqslant 1\)。反之，若 \(|\xi_n| = 1\) 对所有 \(n\) 成立，则 \(x\) 是极点，因为关系式 \(\xi_n = \lambda \eta_n + (1 - \lambda) \zeta_n\) 连同 \(|\eta_n| \leqslant 1\)、\(|\zeta_n| \leqslant 1\) 与 \(0 < \lambda < 1\) 蕴含 \(\xi_n = \eta_n = \zeta_n\)。

\begin{enumerate}
\def\labelenumi{\arabic{enumi})}
\setcounter{enumi}{2}
\tightlist
\item
  设 \(u: \mathrm{E} \to \mathrm{E}'\) 是仿射空间 \(\mathrm{E}\) 到仿射空间 \(\mathrm{E}'\) 中的仿射映射；设 \(\mathrm{C} \subset \mathrm{E}\)、\(\mathrm{C}' \subset \mathrm{E}'\) 是满足 \(u(\mathrm{C}) \subset \mathrm{C}'\) 的两个凸集。若 \(x'\) 是 \(\mathrm{C}'\) 的一个极点，且 \(x\) 是 \(u^{-1}(x') \cap \mathrm{C}\) 的一个极点，则 \(x\) 是 \(\mathrm{C}\) 的极点，这由定义 1 即得。
\end{enumerate}

命题 1. --- 设 B 是 A 的极点所成之集，其中 A 是 Hausdorff 局部凸空间 E 中的非空紧凸集，并设 f 是定义在 A 上且上半连续的凸函数。则 f 在 A 中的上确界可在 B 的（至少）一个点上达到。

用 F 表示由 A 的这样的子集 X 组成的族：X 非空、闭，且包含在 A 中并与 X 相交的每个开线段必含于 X。它具有下列性质；

\begin{enumerate}
\def\labelenumi{(\roman{enumi})}
\item
  A 属于 \(\mathfrak{F}\)。
\item
  点 \(a \in \mathbf{A}\) 使 \(\{a\} \in \mathfrak{F}\)，当且仅当 \(a\) 是 \(\mathbf{A}\) 的一个极点。
\item
  每个非空交 \(\mathbf{X}\)（族 \((\mathbf{X}_{\alpha})\) 的交，该族由 \(\mathfrak{F}\) 中的集组成）也属于 \(\mathfrak{F}\)。
\end{enumerate}

性质 (i)、(ii) 与 (iii) 由定义立得。

\begin{enumerate}
\def\labelenumi{(\roman{enumi})}
\setcounter{enumi}{3}
\tightlist
\item
  设 \(X \in F\)，h 是 A 中凸且上半连续的函数；则 X 中使限制 \(h|X\) 在 X 中达到其上确界的点所成的集 Y 满足：Y 属于 F。
\end{enumerate}

因为，\(h|X\) 在 \(X\) 上上半连续，故其上确界 \(\alpha\)（在 \(X\) 上取的上确界）在 \(X\) 的至少一个点上达到（GT, IV, § 6.2, 定理 3）；于是 \(Y\) 非空，它也是闭的（GT, IV, § 6.2, 命题 1）。另一方面，设 \(x, y\) 是 \(A\) 的两个不同点，并设 \(z = \lambda x + (1 - \lambda)y\) 是 \(Y\) 中满足 \(0 < \lambda < 1\) 的一点；由于 \(Y \subset X\) 且 \(X \in \mathfrak{F}\)，我们有 \(x \in X\) 与 \(y \in X\)；另一方面，由于 \(h\) 是凸的，我们有

\[
h (z) \leqslant \lambda h (x) + (1 - \lambda) h (y)
\]

但由于 \(h(x) \leqslant \alpha\)、\(h(y) \leqslant \alpha\) 且 \(h(z) = \alpha\)，必有 \(h(x) = h(y) = \alpha\)，也就是说 \(x \in \mathbf{Y}\) 且 \(y \in \mathbf{Y}\)。因此 \(\mathbf{Y} \in \mathfrak{F}\)。

在建立了这些性质之后，设 M 是使 f 在 A 中达到其上确界的 \(x \in A\) 所成之集；由 (iv)，\(M \in F\)。另一方面，由 (iii) 以及 F 中的集都是紧集 A 的闭子集这一事实，可知 F 关于序关系 \(\supset\) 是归纳的。由 S, III, §2.4 的定理 2，M 包含一个子集 N，它是 F 的极小元。我们将证明 N 由单独一点组成，从而完成命题的证明。由于 E 是 Hausdorff 局部凸空间，只需证明 E 上每个连续线性型 u 在 N 中是常数（II, p.~38, 推论 1）。而由 (iv) 推出，集 \(N'\)——由 \(x \in N\) 中使 \(u|N\) 在 N 中达到其上确界者组成——满足 \(N'\) 属于 F；由于 N 在 F 中极小，必有 \(N' = N\)。

推论. --- 设 A 是 Hausdorff 局部凸空间 E 中的紧凸集。则 A 的每个闭支撑超平面 H 至少包含 A 的一个极点。

因为，若 \(f(x) = \alpha\) 是 \(H\) 的一个方程，且 \(f(x) \leqslant \alpha\) 在 \(A\) 中成立，只需对 \(f\) 应用命题 1。

定理 1（Krein-Milman）. --- 在 Hausdorff 局部凸空间 E 中，每个紧凸集 A 都是它的极点所成之集的闭凸包。

因为，设 C 是 A 的极点所成之集的闭凸包；显然 \(C \subset A\)。为看出 \(A \subset C\)，只需证明：若 u 是在 E 中连续的仿射线性函数，且在 C 中 \(u(x) \geqslant 0\)，则在 A 中也有 \(u(x) \geqslant 0\)（II, p.~39, 推论 4）；而这由对 -u 应用命题 1 即得。

命题 2. --- 设 x 是 Hausdorff 局部凸空间 E 中紧凸集 A 的一个极点。则对 E 中 x 的每个开邻域 V，存在 E 中的一个开半空间 F，使得 \(x \in F \cap A \subset V \cap A\)（换言之，包含 x 的诸开半空间在 A 上的迹构成 x 在 A 中的基本邻域系）。

对 E 的每个包含 x 的开半空间 D，集 \(A \cap \overline{D}\) 是 x 在 A 中的一个紧邻域，而所有这些邻域的交恰为点 x（任意两个不同的点可被一个闭超平面严格分离（II, p.~38, 命题 4）。由 GT, I, § 9.2 的命题 1，只需证明诸集 \(A \cap \overline{D}\) 构成一个滤子基。现在，若写 \(L_{D} = A \cap (E - D)\)，则集 \(L_{D}\) 是凸的、紧的，且含于凸集 \(A - \{x\}\)；若 \(D_{1}, D_{2}\) 是 E 的两个包含 x 的开半空间，则 \(L_{D_{1}} \cup L_{D_{2}}\) 的凸包 B 因此含于 \(A - \{x\}\)；但 B 是紧集（II, p.~14, 命题 15），故存在一个闭超平面 H 把 x 与 B 严格分离（II, p.~38, 命题 4），而若由 H 决定且包含 x 的开半空间是 D，则我们有 \(L_{D_{1}} \cup L_{D_{2}} \subset L_{D}\)，因此 \(A \cap \overline{D} \subset (A \cap \overline{D}_{1}) \cap (A \cap \overline{D}_{2})\)。

推论. --- 在 Hausdorff 局部凸空间中，设 K 是紧凸集 A 的一个紧子集。则下列条件等价。

\begin{enumerate}
\def\labelenumi{\alph{enumi})}
\item
  A 是 K 的闭凸包。
\item
  K 与每一集相交，该集是 A 与其某一支撑超平面的交。
\item
  K 包含 A 的极点所成之集。
\end{enumerate}

\(a) \Rightarrow b)\)。假设存在 A 的一个支撑超平面 H，其方程为 \(f(x) = \alpha\)，使得 \((H \cap A) \cap K = \varnothing\)，并例如假设在 A 中 \(f(x) \geqslant \alpha\)。由于按假设 \(f(x) - \alpha > 0\) 对所有 \(x \in K\) 成立，且 K 是紧的，我们有

\[
\beta = \inf _ {x \in \mathbf {K}} f (x) > \alpha ,
\]

于是 K 含于闭半空间 \(f(x) \geqslant \beta\)；因此对 K 的闭凸包 A 同样如此，而这是荒谬的。

\(b) \Rightarrow c)\)。假设 \(x\) 是 \(\mathbf{A}\) 的一个不属于 \(\mathbf{K}\) 的极点；则存在 \(\mathbf{V}\)（\(x\) 在 \(\mathbf{E}\) 中的一个邻域），使得 \(\mathbf{V} \cap \mathbf{A} \cap \mathbf{K} = \varnothing\)。但由命题 2，我们可以假设 \(\mathbf{V}\) 是由超平面 \(\mathbf{H}\)（其方程为 \(f(z) = \alpha\)）所定义的开半空间。若例如 \(f(x) > \alpha\)，则对所有 \(y \in \mathbf{K}\)，我们有 \(f(y) \leqslant \alpha\)，因此 \(\mathbf{K}\) 不与 \(\mathbf{A}\) 和支撑超平面 \(f(z) = \gamma > \alpha\)（平行于 \(\mathbf{H}\)）的交相遇（II, p.~37, 命题 2）；这是荒谬的。

\(c) \Rightarrow a)\)。这是 Krein-Milman 定理的明显推论。

习题 11）。定理 1（II, p.~56）的证明表明，无论如何，A 都是 A 中属于某一支撑超平面的那些极点所成之集的闭凸包。

